\clearpage
\SourceLeaf{407}\pdfbookmark[1]{III. Zweite Methode. Mai–Juli 1805}{ceres-second-method}
\begin{center}
{\large BRIEFWECHSEL UND NACHLASS.\par}
\bigskip
[III. ZWEITE METHODE.]\par
[MAI—JULI 1805.]\par
\bigskip
[1.]\par
\medskip
\textsc{Gauss} an \textsc{Olbers}, Braunschweig, 10. Mai 1805.
\end{center}
\noindent\hspace*{\parindent}.\quad.\quad.\quad.\quad.\quad Die Methode, nach der ich die {\AstroSymbols ⚳}-Störungen zu berechnen angefangen hatte, habe ich doch wieder aufgegeben. Die allzu vielen eintönigen mechanischen Rechnungen, die ich dabei vor mir sah, haben mich abgeschreckt. Selbst wenn alle Rechnungen, die ich anderen hätte übertragen können, von Herrn \textsc{Bessel} und Herrn von \textsc{Lindenau}\quad.\quad.\quad.\quad.\quad übernommen worden wären, wäre für mich noch mehr übrig geblieben, als meine Geduld hätte bewältigen können.

Ich habe indessen bereits eine andere Methode ersonnen, die ebenso weit führen kann wie jene, aber bei weitem weniger, wenn auch anspruchsvollere Arbeit erfordert. Ich bin bei ihrer Anwendung auf Ceres schon ein gutes Stück vorangekommen, vorerst allerdings nach einem eingeschränkteren Plan, da ich nur bis zu den 5. Potenzen der Exzentrizitäten von {\AstroSymbols ♃} und {\AstroSymbols ⚳} gehe. Diese Methode hat für mich umso mehr Reiz, als ich dabei viele, schon vor längerer Zeit angestellte, ziemlich tiefgehende Untersuchungen über besondere Arten transzendenter Funktionen gut verwenden kann. Ich werde versuchen, Ihnen in [der] Folge eine Vorstellung davon zu geben. Ich hoffe auch, alles so einrichten zu können, dass mir Unterstützung von anderer Seite die Arbeit erheblich erleichtert. Für meine diesmalige Rechnung für {\AstroSymbols ⚳} kommt das zu spät (denn gerade in dem Teil, in dem Hilfe nötig wäre, bin ich selbst schon zu weit \SourceLeaf{408}vorgerückt). Wenn ich sie wiederhole, wird es aber möglich sein. Die Wiederholung wird nötig sein, da die erweiterten Störungsgleichungen zweifellos noch mit erheblichen Änderungen der Elemente selbst verbunden sein werden. Ebenso wäre Unterstützung möglich, wenn ich diese Arbeit einmal für {\AstroSymbols ⚴} und {\AstroSymbols ⚵} vornehme, wo sie beträchtlich umfangreicher sein wird.\quad.\quad.\quad.\quad.\quad.
\par\medskip
\begin{center}\rule{.30\linewidth}{.4pt}\end{center}
